%- block legend

\begin{center}
    \includegraphics[scale=1.0]{vaca.png}
\end{center}

Quando começou a estudar números binários, a vaca MU ficou fascinada com o
assunto!

Ela escolheu dois números inteiros positivos $A$ e $B$, com $A \leq B$,
e escreveu a representação binária, sem zeros à esquerda, de todos os números de $A$ até $B$ (inclusive
e nesta ordem) em seu quadro branco, formando uma string $S_{A,B}$. Depois, ela
contou quantos dígitos 1 (um) apareceram em $S_{A,B}$ e anotou este número no
quadro, o chamando de $f(A,B)$.

Como exemplo, considere $A=5$ e $B=10$. A string $S_{5,10}$ é
101110111100010011010. Como há 12 dígitos 1 (um) em
$S_{5,10}$, tem-se que $f(5,10) = 12$.

O camaleão Kurt resolveu pregar uma peça na vaca MU e, durante a noite, apagou
várias informações que estavam no quadro! Os únicos valores que o camaleão
deixou escritos são $B$ e $f(A,B)$.

Sua tarefa é ajudar a vaca MU a lembrar qual era o valor de $A$.

%- endblock

%- block input
A primeira linha contém um inteiro $T$ ($1 \leq T \leq 1000$), o número de casos
de teste.

Cada caso de teste é dado por uma linha contendo dois inteiros $B$ e $f$
($1 \leq B \leq 10^{16}$, $0 \leq f \leq 10^{18}$), o valor de $B$ e de
$f(A,B)$, respectivamente.

%- endblock

%- block output
Para cada caso de teste, imprima uma linha com o valor de $A$.

Caso não exista um valor de $A$ condizente com os valores de $B$ e $f(A,B)$
dados na entrada, imprima a mensagem dada no exemplo de
saída abaixo. Note que a mensagem deve ser impressa em \textit{três} linhas.
%- endblock

%- block editorial
Seja $a$ um inteiro positivo menor ou igual a $B$. Note que:
\begin{itemize}
    \item se $f(a,B) = f(A,B)$, então $A=a$;
    \item se $f(a,B) < f(A,B)$, então $A<a$; 
    \item se $f(a,B) > f(A,B)$, então $A>a$.
\end{itemize}

Esta observação nos permite resolver o problema fazendo \textit{busca binária} em $a$.
Assim, o problema se torna calcular o valor de $f(a,B)$ para
algum $a \leq B$ dado.

Seja $F(N)$ o número total de \textit{bits} 1 em todos os inteiros
menores ou iguais a $N$. Note que
\begin{itemize}
    \item $f(a,B) = F(B) - F(a-1)$
\end{itemize}

Assim, o problema agora se reduz a calcular $F(N)$ para algum inteiro $N$ dado.

$F(N)$ pode ser computado com uma técnica de Programação Dinâmica conhecida como
\textit{digit dp}.

Seja $n$ o número de \textit{bits} de $N$, e seja $N_{n-1}N_{n-2}...N_1N_0$ a
representação binária de $N$, do \textit{bit} mais para o menos significativo.

Considere que você está escrevendo, em binário, os números menores ou iguais a
$N$, do \textit{bit} mais significativo ao menos significativo. Considere o estado
$dp[b][l]$, onde $b$ indica a posição do \textit{bit} que você está escrevendo
agora, e
$l=1$ (de
``limitado'') se todos os \textit{bits} que você já escreveu até o
momento são exatamente iguais a $N_{n-1}N_{n-2}...N_{b+1}$, ou $l=0$ caso contrário
(a ideia aqui é que, se $l=1$, você ainda precisa ``tomar cuidado'' para não
 escrever um número maior que $N$; se $l=0$, você já pode escrever qualquer
 número, pois ele certamente será menor ou igual a $N$. O ignorar do último caso
 da $dp$, dado abaixo, garante esta propriedade).
O estado $dp[b][l]$ retorna o número de \textit{bits} 1 que serão escritos ao
todo sob estas condições, de forma que $F(N) = dp[n-1][1]$.

O valor de $dp[b][l]$ é:
\begin{itemize}
    \item considerando escrever o \textit{bit} 0 na posição $b$:
    \begin{itemize}
        \item[$\rightarrow$] $dp[b-1][l']$, onde $l' = 1$ se e somente se $l = 1$ e $N_b = 0$;
    \end{itemize}
    \item considerando escrever o \textit{bit} 1 na posição $b$:
    \begin{itemize}
        \item[$\rightarrow$] se $l=0$, o \textit{bit} 1 aparece, na posição $b$, em todos os $2^b$ números
        possíveis de serem escritos a partir desse estado; logo, somar
        $2^b + dp[b-1][0]$;
        \item[$\rightarrow$] se $l=1$ e $N_b = 1$, o \textit{bit} 1 aparece, na posição $b$, em
        todos os $N_{b-1}N_{b-2}...N_0 + 1$ números possíveis de serem escritos a
        partir desse estado; logo, somar
        $N_{b-1}N_{b-2}...N_0 + 1 + dp[b-1][1]$;
        \item[$\rightarrow$] se $l=1$ e $N_b = 0$, não é possível escrever o \textit{bit} 1 na
        posição $b$. Ignore este caso.
    \end{itemize}
\end{itemize}

Uma possível referência para \textit{digit dp} é \url{https://www.geeksforgeeks.org/dsa/digit-dp-introduction/}

Complexidade: $O(\lg B)$ (da busca binária) $\times$ $O(\lg B)$ (do
        \textit{digit dp}) = $O((\lg B)^2)$ por caso de teste.

%- endblock
